\pagestyle{myheadings}
\chapter{MARCO TEÓRICO REFERENCIAL}
\section{Antecedentes}

\subsection{Antecedentes internacionales}

~\cite{Aslam2019} realizó un estudio titulado:Design Changes in Construction Projects – Causes and Impacton the Cos, en la ciudad de Islamabad,Pakistan.

El objetivo general del trabajo fue examinar el impacto de los cambios de diseño en el coste del proyecto e identificar las acciones responsables de estos cambios. La muestra estuvo constituida por .......,el diseño que utilizó fue la revisión sistemática de la literatura anterior publicada en revistas bien establecidas, y se analizaron los contenidos.Los instrumentos que usó fue .........y los resultados obtenidos han sido que se estableció que el cambio de diseño es uno de los factores predominantes para el sobrecoste, y en algunos casos, puede dar lugar a un sobrecoste de entre el 5 y el 40\% del coste del proyecto. 

~\cite{Mansur2019} realizó un estudio titulado:Components of Cost Overrun in China Construction Projects en la ciudad de  Skudai, Malaysia.

El objetivo de este estudio fue evaluar los principales factores que influyen en los componentes de los costes, como la mano de obra, los materiales, los equipos y los subcontratistas.La muestra estuvo constituida por .......El diseño que se utilizó fue ........y los instrumentos que se usaron fueron 2 encuestas ........los resultados que obtuvieron fue que Los resultados de la encuesta indicaron que la baja productividad de la mano de obra, la escalada de los precios de los materiales, el elevado coste de la maquinaria y el flujo de caja, y la mala gestión financiera son las posibles causas relacionadas con el exceso de costes en los proyectos de construcción de China, en comparación con otros factores.

~\cite{Susanti2020} realizó un estudo titulado:\textit{Cost Overrun in Construction Projects in Indonesia} en la ciudad de Java Central en Indonesia.
El objetivo general del estudio fue identificar los factores causantes de los sobrecostes en un proyecto y analizar los factores que más influyen en la aparición de sobrecostes en proyectos de construcción en Indonesia, tal y como los perciben el propietario y el contratista.La muestra estuvo constituida por ............El diseño que se utilizó fue .................Los instrumentos que se usaron fue un cuestionario y los resultados obtenidos han sido que se identificaron 15 factores: retraso en la disponibilidad de la obra; condiciones de la obra; condiciones sociales de la obra; orden de cambio; retrabajo; rendimiento del subcontratista y/o proveedor; retraso en la aprobación/permiso; inexactitud en la presupuestación; programación y planificación de recursos; fluctuaciones en el precio de los materiales; normas y reglamentos; requisitos adicionales del propietario; inflación; retraso en el pago; escasa liquidez y mal tiempo. 

~\cite{Plebankiewicz2018} realizó un estudio titulado: \textit{Model of Predicting Cost Overrun in Construction Projects} en la ciudad de Kraków,Polonia.El objetivo general del estudio fue describir el concepto de un modelo de predicción del aumento de los costes de las obras de construcción en relación con los previstos.La muestra estuvo constituida por ..................El diseño que se utilizó fue..............., los instrumentos que se utilizaron fue............Los instrumentos que utilizaron fue...................y los resultados obtenidos fueron que ..............

\subsection{Antecedentes nacionales}

~\cite{FloresBedon2022} realizó un estudio titulado:Evaluación de la variación presupuestal en las obras ejecutadas por impuestos y por contrata, en el Gobierno Regional Ucayali, período 2019.El objetivo general del estudio fue describir el concepto de un modelo de predicción del aumento de los costes de las obras de construcción en relación con los previstos.La muestra estuvo constituida por 17 obras ejecutadas por contrata e impuestos.El diseño que se utilizó fue correlacional, cuantitativo y no experimental,Los instrumentos que utilizaron fue fichas de observación y los resultados obtenidos fueron que que el costo de los adicionales varia en un 32\% respecto al presupuesto inicial.


~\cite{FigueroaVillarreal2022} realizó un estudio titulado:Gestión de riesgos en la planificación de la ejecución de obra de la institución educativa Santo Tomás de Cocharca, Oyón,2021.describir los procesos de la gestión de riesgos en la planificación de la ejecución de obra de la Institución Educativa Santo Tomás de Cochamarca, Oyón, 2021.Método: el tipo de investigación es básica, de enfoque cuantitativo-cualitativo (mixto),nivel descriptivo y diseño no experimental. Se analizó el expediente técnico de código2339464, se realizó la recolección de información y su respectivo análisis, la lista de riesgos y los anexos según la directiva. Luego, se analizó e interpretó la información vinculada a los 10 riesgos identificados, mediante la estadística descriptiva empleando el programa Excel.

~\cite{MachadoYnquillay2021} realizó un estudio titulado:Análisis de riesgos para la determinación de la variación costo – tiempo en la planificación de la obra''construcción integral sistema de riego Pampaconga – Limatambo – Anta – Cusco'' ejecutado por el Proyecto Especial Sierra Centro Sur.El objetivo general fue determinar los riesgos que influyeron en Costo y Tiempo.La muestra estuvo constituida por los WPS y EPS de la obra en mención.El diseño que se utilizó fue mixto, los instrumentos que utilizó fue instrumentos documentales de observación: Expediente técnico, Cuadernos de obra e informes técnicos mensuales. También utilizaron instrumentos elaborados de comunicación:listas de verificación, tormenta de ideas, encuestas estructuradas y los resultados obtenidos fue variación en el Costo del 12.53\% con respecto al presupuesto programado y una variación en el Tiempo del 68.81 \% con respecto al tiempo programado.

~\cite{FontanaBecerra2022} realizó un estudio titulado: Las ampliaciones de plazo en obras públicas y su relación con la penalidad por mora en la Municipalidad Distrital de Cerro Colorado.El objetivo general del estudio fue determinar la relación entre las ampliaciones de plazo y la penalidad por mora en las obras públicas en la Municipalidad Distrital de Cerro Colorado-Arequipa, 2020-2021.La muestra estuvo constituida por  11 expedientes de ampliación de plazo, derivadas de 06 obras públicas de la Municipalidad Distrital de Cerro Colorado, durante los años 2020-2021.El diseño que se utilizó fue descriptivo.Los instrumentos que utilizaron fue Ficha de observación y los resultados obtenidos fueron incidencia de aprobación de las ampliaciones es de 81.8 \%, influyendo en la no aplicación de la penalidad por mora, sustentándose por el atraso en las obras, en arreglo a la ley.

\subsection{Antecedentes regionales o locales}

No se encontraron bibliografías concernientes al tema de investigación en los repositorios de las 2 únicas universidades de la zona de estudio.